\section{Related Work}\label{sec:related}

\textbf{Certifiable layers and the SDP route.} SDPRLayers \cite{holmes2025sdpr} lifts a homogeneous QCQP to $X\succeq0$, solves the Shor relaxation, and differentiates the original KKT system with an implicit-selection theorem that tolerates redundant constraints. Its released pose estimator uses a $13$-dimensional vector $x=[h;\mathrm{vec}(C);t]$ and 24 quadratic constraints (6 orthogonality, 9 handedness, 9 norm), with the last 18 marked redundant. Our work differs in four ways. (i) \emph{Variable:} we optimise over $q\in\mathbb{S}^3$ (four variables, one constraint) and eliminate $t$ exactly, so the dimension of the nonconvex problem is constant and independent of the number of correspondences. (ii) \emph{Global mechanism:} QPEP enumerates all KKT roots and compares objectives; there is no convex relaxation and hence no tightness assumption. (iii) \emph{Backward:} we differentiate the reduced stationarity system with a $3\times3$ Riemannian Hessian; SDPRLayers differentiates a KKT system with redundant rows after row reduction. (iv) \emph{Ambiguity:} we return a smoothed posterior over modes; SDPRLayers returns a single certifiable solution and has no defined behaviour when the relaxation is not tight. We do \emph{not} claim that our $10\times10$ analysis applies to the released 24-constraint implementation, whose constraint span is 17-dimensional; the full-20D $13\times13$ relaxation, to which it does apply, is proved equivalent to the $10\times10$ one in Theorem~\ref{thm:thirteen}.

\textbf{Probabilistic pose estimation.} EPro-PnP \cite{chen2022epropnp} defines a Boltzmann distribution on $SE(3)$ and trains with a KL loss using adaptive multiple importance sampling. Our posterior is also Boltzmann, but the support is $\mathbb{S}^3/\{\pm1\}$ rather than $SE(3)$, translation is marginalised exactly, and the rotation integral is approximated analytically over an exhaustive list of modes rather than sampled. The trade is bias for variance: EPro-PnP is asymptotically exact but high-variance, whereas our layer is deterministic but has an $O(\tau)$ Laplace bias away from bifurcations. A fair comparison must separate these two error sources, and Section~\ref{sec:experiments} does not claim to have made it: only the discrete-mode sampling variance is measured.

\textbf{The QPEP solver.} Wu \emph{et al.} \cite{wu2022qpep} unify a large class of pose problems as quadratic problems in a unit quaternion, eliminate translation, solve the resulting polynomial system with a Gr\"obner basis, and analyse solvability, observability and covariance. Their solver returns at most 40 candidate complex roots in the released implementation (the generic 40-root count is confirmed by a total-degree homotopy in Section~\ref{sec:rootcount}); the global solution is obtained by comparing objectives, not by a convex certificate. QPEP provides no derivative with respect to the data. Our contribution is precisely the missing differentiable layer: the closed-form Euler multiplier, the $3\times3$ implicit differentiation, the antipodal invariance, and the Laplace posterior over the enumerated roots. We rely on QPEP for the forward enumeration and treat its completeness as an explicit assumption (A2).

\textbf{Implicit differentiation of local solvers.} BPnP \cite{chen2020bpn} backpropagates through a PnP solver by implicit differentiation of its optimality conditions. The solver is local and initialised from a network prediction; the gradient is that of the local solution, not of the global one. At a branch crossing the local solution may be the wrong one, and the derivative of the global argmin does not exist as a classical function. Theseus \cite{pineda2022theseus} provides a differentiable nonlinear least-squares library with Lie-group support, sparse solvers, batching and implicit/direct loss minimisation; it is a general-purpose local optimiser, does not enumerate global modes, and provides no certificate. Our layer is narrower (quartic costs in a quaternion) but exhaustive, with a closed-form backward pass of constant size.

\textbf{Polynomial optimisation, moments and sums of squares.} The relaxation we analyse is the degree-two Shor/Veronese relaxation, not a full moment hierarchy. Lasserre's hierarchy \cite{lasserre2001} and Laurent's survey \cite{laurent2009} provide the general moment/SOS framework; Hilbert \cite{hilbert1888} established that nonnegative forms need not be sums of squares, and Blekherman \cite{blekherman2006} quantifies how many more nonnegative forms there are than SOS forms. The Choi--Lam quartic \cite{choi1977extremal,choi1980zeros} is the standard explicit witness of that gap, which is why we use it to construct an instance with a certified relaxation gap; the algebraic geometry of the Veronese variety behind Theorem~\ref{thm:lift} is standard \cite{harris1992}. We emphasise what we do \emph{not} identify: our SDP is not the textbook second-order Lasserre moment relaxation with a sphere localising matrix, and we make no claim about higher-order hierarchies.

\textbf{Laplace asymptotics and the geometry of $SO(3)$.} The Laplace approximation and its uniform (Airy) counterpart are classical asymptotic-expansion results \cite{bleistein1986,wong1989}; the branch-merging analysis of Section~\ref{sec:discussion} is exactly the degenerate case in which the standard expansion ceases to be uniform. The measure constants in the rotation integral come from the Haar measure of $SO(3)$ and the double cover $\mathbb{S}^3\to SO(3)$ \cite{chirikjian2012}. The implicit-function theory behind Theorems~\ref{thm:implicit}--\ref{thm:cert} is that of Dontchev and Rockafellar \cite{dontchev2009}, and the numerics of the forward enumeration (damped Newton, Levenberg--Marquardt, KKT systems) is standard \cite{nocedal2006}.

\textbf{Scope relative to companion tracks.} The companion track \textsc{Constrained-QPEPLayer} adds inequality constraints (handedness, non-penetration, frustum) through an augmented Lagrangian and a Fischer--Burmeister reduction; \textsc{Robust-QPEPLayer} adds outlier-robust kernels and the curvature correction of the IRLS weights. The present layer assumes neither: its $M(\theta)$ may be any weighted Gram matrix, and its theory does not depend on unpublished results from the other tracks.
